\documentclass[aspectratio=169,12pt]{beamer}

\usepackage{ctex}
\usepackage{amsmath}
\usepackage{amssymb}
\usepackage{graphicx}
\usepackage{booktabs}
\usepackage{colortbl}
\usetheme{Madrid}

\title{矩阵对角化与若当标准型}
\subtitle{从特征向量到若当块}
\author{}
\date{2026年3月}

\begin{document}

\frame{\titlepage}

\section{前置知识}

\begin{frame}
\fontsize{11}{13}\selectfont
\centerline{\textbf{特征值与特征向量的几何意义}}

\bigskip
\centerline{\textbf{问题}：什么是线性变换的「不变方向」？}

\bigskip
\pause
如果向量 $v \neq 0$ 满足 $Av = \lambda v$，则：
\begin{itemize}
    \item $v$ 是特征向量
    \item $\lambda$ 是特征值
\end{itemize}

\pause
\bigskip
特征向量 = 变换后方向不变的向量

特征值 = 在该方向上的伸缩倍数
\end{frame}

\begin{frame}
\fontsize{10}{12}\selectfont
\centerline{\textbf{特征值的几何意义}}

\[
\lambda > 1 : \text{拉伸}
\]
\[
0 < \lambda < 1 : \text{压缩}
\]
\[
\lambda = 1 : \text{方向完全不变}
\]
\[
\lambda < 0 : \text{反向并伸缩}
\]
\end{frame}

\section{问题引入}

\begin{frame}
\fontsize{11}{13}\selectfont
\centerline{\textbf{我们为什么要研究这个问题？}}

\bigskip
\centerline{在研究矩阵时，我们希望把矩阵变得更简单。}

\bigskip
\pause
对角矩阵是最简单的形式：
\[
\Lambda = \begin{pmatrix}
\lambda_1 & & \\
& \ddots & \\
& & \lambda_n
\end{pmatrix}
\]

\pause
\bigskip
幂运算非常容易：
\[
\Lambda^k = \begin{pmatrix}
\lambda_1^k & & \\
& \ddots & \\
& & \lambda_n^k
\end{pmatrix}
\]

\pause
\bigskip
\centerline{\textbf{核心问题}：是不是所有矩阵都能变成对角矩阵？}
\end{frame}

\section{对角化}

\begin{frame}
\fontsize{11}{13}\selectfont
\centerline{\textbf{矩阵对角化的定义}}

\bigskip
若存在可逆矩阵 $P$ 和对角矩阵 $\Lambda$，使得
\[
A = P\Lambda P^{-1}
\]

则称矩阵 $A$ 可对角化，$P$ 的列是 $A$ 的特征向量。
\end{frame}

\begin{frame}
\fontsize{10}{12}\selectfont
\centerline{\textbf{充要条件}}

\bigskip
$A$ 可对角化的充分必要条件：

\bigskip
\centerline{\textbf{$A$ 有 $n$ 个线性无关的特征向量}}

\bigskip
\pause
几何重数 = 代数重数
\end{frame}

\section{不能对角化}

\begin{frame}
\fontsize{10}{12}\selectfont
\centerline{\textbf{不能对角化的情况}}

\bigskip
考虑矩阵 $A = \begin{pmatrix} 2 & 1 \\ 0 & 2 \end{pmatrix}$

\pause
\bigskip
特征值 $\lambda = 2$（重根），只有1个特征向量！

\pause
\bigskip
\begin{center}
\begin{tabular}{|c|c|c|}
\hline
概念 & 含义 & 本例 \\
\hline
代数重数 & 特征值作为方程根的重数 & 2 \\
\hline
几何重数 & 线性无关特征向量的个数 & 1 \\
\hline
\end{tabular}
\end{center}

\pause
\bigskip
几何重数 < 代数重数 → 不能对角化
\end{frame}

\section{若当块的发现}

\begin{frame}
\fontsize{11}{13}\selectfont
\centerline{\textbf{核心问题}：如何在特征值旁边「添加」非零元素？}

\bigskip
我们想要：
\[
A^k = P J^k P^{-1}
\]

其中 $J$ 是「类似对角」的矩阵。

\pause
\bigskip
如果不能完全对角化，能否在紧邻对角线的位置添加非零元素？
\[
\begin{pmatrix}
\lambda & ? \\
0 & \lambda
\end{pmatrix}
\]
\end{frame}

\begin{frame}
\fontsize{10}{12}\selectfont
\centerline{\textbf{填什么最简单？}}

\bigskip
选项1：填 $\lambda$ 本身？$\begin{pmatrix} \lambda & \lambda \\ 0 & \lambda \end{pmatrix}$

\bigskip
选项2：填1？$\begin{pmatrix} \lambda & 1 \\ 0 & \lambda \end{pmatrix}$

\pause
\bigskip
\centerline{\textbf{填1最简单！}}

\begin{itemize}
    \item 任何非零数都可以通过「缩放」变成1
    \item 幂运算时，1的幂永远是1，形式最简洁
    \item 上移位矩阵 $N$ 有性质：$N^k$ 形式规则
\end{itemize}
\end{frame}

\begin{frame}
\fontsize{10}{12}\selectfont
\centerline{\textbf{验证}}

\[
\begin{pmatrix} \lambda & 1 \\ 0 & \lambda \end{pmatrix}^2
= \begin{pmatrix} \lambda^2 & 2\lambda \\ 0 & \lambda^2 \end{pmatrix}
\]

\bigskip
\pause
这就是\textbf{若当块}！
\end{frame}

\section{广义特征向量}

\begin{frame}
\fontsize{10}{12}\selectfont
\centerline{\textbf{第一层思考}：只有1个特征向量，怎么办？}

\bigskip
已知：特征向量满足 $(A-\lambda I)v_1 = 0$

\bigskip
问题：我有2维空间，但只有1个特征向量，如何找到另一个「有用」的方向？
\end{frame}

\begin{frame}
\fontsize{10}{12}\selectfont
\centerline{\textbf{第二层思考}：能否放松条件？}

\bigskip
特征向量要求：变换后完全不变（方向不变）

\bigskip
能否降低要求：变换后「差不多」不变？

\bigskip
\pause
类比：如果找不到完美的，找「次优」的也行！
\end{frame}

\begin{frame}
\fontsize{10}{12}\selectfont
\centerline{\textbf{第三层思考}：如何描述「差不多」？}

\bigskip
思路：递推——既然 $v_1$ 是不变的，能否通过 $v_1$ 找到下一个？
\[
(A-\lambda I)v_2 = v_1
\]

\pause
\bigskip
为什么这样想？
\begin{itemize}
    \item 如果 $v_1$ 是「一次变换后不变」
    \item 那么 $v_2$「两次变换后」应该也不变
    \item 验证：$(A-\lambda I)^2 v_2 = (A-\lambda I)v_1 = 0$
\end{itemize}
\end{frame}

\begin{frame}
\fontsize{10}{12}\selectfont
\centerline{\textbf{第四层思考}：链条的源头是特征值}

\bigskip
\centerline{\textbf{核心发现}：所有链条都源于特征值！}

\bigskip
每一代的关系：
\begin{itemize}
    \item $J v_1 = \lambda v_1$（特征向量，紧邻特征值）
    \item $J v_2 = \lambda v_2 + v_1$（广义特征向量，指向 $v_1$）
    \item $J v_3 = \lambda v_3 + v_2$（广义特征向量，指向 $v_2$）
\end{itemize}

\pause
\bigskip
这就是链条！源头是特征值 $\lambda$，每一代「指向」前一代。
\end{frame}

\section{若当块定义}

\begin{frame}
\fontsize{11}{13}\selectfont
\centerline{\textbf{若当块定义}}

\[
J_k(\lambda) = 
\begin{pmatrix}
\lambda & 1 & & \\
& \lambda & \ddots & \\
& & \ddots & 1 \\
& & & \lambda
\end{pmatrix}_{k \times k}
\]

\bigskip
\begin{itemize}
    \item 对角线：$\lambda$（特征值，链条源头）
    \item 上对角线：1（链条传递）
    \item 几何：每一代「指向」前一代
\end{itemize}
\end{frame}

\section{若当块结构判断}

\begin{frame}
\fontsize{10}{12}\selectfont
\centerline{\textbf{新问题}：4阶矩阵的情况}

\bigskip
如果是4阶矩阵，特征值 $\lambda = 2$（四重），但只有2个特征向量，会有什么情况？

\bigskip
两种可能：
\begin{enumerate}
    \item 两个2阶块：$J_2(2) + J_2(2)$
    \item 一个3阶块 + 一个1阶块：$J_3(2) + J_1(2)$
\end{enumerate}

\pause
\bigskip
\centerline{\textbf{如何判断是哪种？}}
\end{frame}

\begin{frame}
\fontsize{10}{12}\selectfont
\centerline{\textbf{方法：通过秩判断}}

\bigskip
设 $N = A - \lambda I$，计算 $N^k$ 的秩。

\begin{center}
\begin{tabular}{|c|c|c|c|}
\hline
结构 & rank(N) & rank(N²) & rank(N³) \\
\hline
两个2阶块 & 2 & 0 & 0 \\
\hline
一个3阶+一个1阶 & 2 & 1 & 0 \\
\hline
一个4阶块 & 3 & 2 & 1 \\
\hline
可对角化 & 2 & 0 & 0 \\
\hline
\end{tabular}
\end{center}

\pause
\bigskip
\centerline{\textbf{规律}：计算 $N^k$ 的秩，观察每次下降多少，就能判断链条的长度分布！}
\end{frame}

\section{若当化定理}

\begin{frame}
\fontsize{11}{13}\selectfont
\centerline{\textbf{若当化定理}}

\bigskip
任意 $n$ 阶矩阵 $A$，都存在可逆矩阵 $P$，使得
\[
A = PJP^{-1}
\]

其中 $J$ 是若当矩阵，由若干若当块组成：
\[
J = \text{diag}(J_{k_1}(\lambda_1), J_{k_2}(\lambda_2), ...)
\]
\end{frame}

\section{统一理解}

\begin{frame}
\fontsize{10}{12}\selectfont
\centerline{\textbf{统一理解：对角化是若当化的特例}}

\begin{center}
\begin{tabular}{|c|c|c|}
\hline
情况 & 若当块结构 & 关系 \\
\hline
可对角化 & $n$个1阶块 & 若当化的特例 \\
\hline
不能对角化 & 若干高阶块 & 一般情况 \\
\hline
\end{tabular}
\end{center}

\pause
\bigskip
\centerline{\textbf{统一理论}：$A = PJP^{-1}$}

\begin{itemize}
    \item $J$ 是对角矩阵 → 对角化
    \item $J$ 是若当矩阵 → 若当化
\end{itemize}

\pause
\bigskip
这就是矩阵分解的完整理论！
\end{frame}

\section{结论}

\begin{frame}
\centerline{\textbf{结论}}

\bigskip
\begin{enumerate}
    \item \textbf{问题}：能否把矩阵变得像对角矩阵一样简单？
    \item \textbf{目标}：对角化（最简单）
    \item \textbf{困难}：几何重数 < 代数重数
    \item \textbf{解决}：若当化（退而求其次）
    \item \textbf{探索}：广义特征向量链
    \item \textbf{方法}：通过 $N = A - \lambda I$ 的秩判断结构
\end{enumerate}
\end{frame}

\end{document}
